%==============================================================================
%  ANNEXE - Principes SOLID et design patterns
%  Application au systeme Zümm : justification "avant / apres"
%  A inserer via %==============================================================================
%  ANNEXE - Principes SOLID et design patterns
%  Application au systeme Zümm : justification "avant / apres"
%  A inserer via %==============================================================================
%  ANNEXE - Principes SOLID et design patterns
%  Application au systeme Zümm : justification "avant / apres"
%  A inserer via %==============================================================================
%  ANNEXE - Principes SOLID et design patterns
%  Application au systeme Zümm : justification "avant / apres"
%  A inserer via \input{annexe_solid_patterns} dans cahier_des_charges.tex
%  (utilise la charte "Miel & noir" : miel, mielfonce, encadre, colonne C,
%   longtable, booktabs -- deja definis dans le preambule principal)
%==============================================================================

%==============================================================================
\chapter{Annexe - Principes SOLID et design patterns}
\label{chap:solid}
%==============================================================================
Cette annexe repond a la question du \emph{comment} de la qualite logicielle.
Elle documente l'application des cinq principes \textbf{SOLID} et d'un ensemble
de \textbf{design patterns} (patrons de conception) au systeme \textbf{Z\"umm},
en mettant chaque decision en regard de la conception initiale (\emph{avant})
et de la conception refactoree (\emph{apres}). L'objectif est une architecture
\textbf{multi-niveaux, evolutive et faiblement couplee}, conforme aux exigences
du chapitre~\ref{chap:conception}.

\begin{encadre}
\textbf{Sources methodologiques :} les definitions retenues suivent la
presentation des principes SOLID de \emph{alexsoyes.com} et le catalogue de
patrons de \emph{refactoring.guru}, adaptes au domaine apicole du projet.
\end{encadre}

%------------------------------------------------------------------------------
\section{Les cinq principes SOLID}
%------------------------------------------------------------------------------
\begin{longtable}{>{\bfseries}p{2.4cm} p{6cm} p{6.2cm}}
\toprule
\textbf{Principe} & \textbf{Avant (conception initiale)} & \textbf{Apres (refactoree)} \\
\midrule
\endhead
S -- Responsabilite unique &
L'entite \texttt{Ruche} porte la methode de calcul
\texttt{getBeeHouseHoneyActualQuantity()} : elle melange l'etat (donnees) et la
logique metier ; \texttt{Visite} regroupe donnees et champs de rapport. &
Le calcul est extrait dans \texttt{ServiceMielImpl} (via le Composite) et le
rapport devient \texttt{RapportVisite}. Chaque classe n'a plus qu'une seule
raison de changer. \\
O -- Ouvert / ferme &
La detection d'alertes et la conversion d'unites reposeraient sur des
\texttt{if / switch} sur \texttt{typeSeuil} et \texttt{unite} : ajouter une
regle ou une unite impose de modifier le code existant. &
\texttt{RegleAlerte} (Strategy) et \texttt{ConvertisseurUnite} : une nouvelle
regle ou unite = une nouvelle classe, sans toucher au \texttt{GestionnaireAlertes}.
Ouvert a l'extension, ferme a la modification. \\
L -- Substitution de Liskov &
Le diagramme de cas d'utilisation generalise l'agent superviseur a partir de
l'agent apiculteur ; transpose tel quel en heritage de classes, un sous-type
restrictif casserait la substituabilite. &
\texttt{Corps} et \texttt{Hausse} heritent de \texttt{Compartiment} en
respectant integralement son contrat (capacite 1..10, \texttt{getPoidsMiel}).
Le role de l'agent passe par l'enumeration \texttt{Role} + permissions
(composition) plutot qu'un heritage restrictif. \\
I -- Segregation des interfaces &
Un service unique \og fourre-tout \fg{} (CRUD + visites + tableaux de bord +
API miel) obligerait l'application tierce a dependre de methodes inutilisees. &
Interfaces ciblees : \texttt{ServiceRequeteMiel} (l'API tierce ne voit que
\texttt{getBeeHouseHoneyActualQuantity}), \texttt{ServiceVisite},
\texttt{ServiceTableauDeBord}. Chaque client ne depend que de son contrat. \\
D -- Inversion de dependance &
La couche metier (EJB) appelle directement JDBC (implementation concrete) :
couplage fort a la base et aux \texttt{ResultSet}. &
La couche metier depend d'abstractions \texttt{RucheRepository} /
\texttt{VisiteRepository} (pattern DAO), implementees par
\texttt{RucheRepositoryJdbc}. Tests par mock et migration JDBC \textrightarrow{}
Spring Data JPA facilites (cf. annexe~\ref{chap:migration}). \\
\bottomrule
\end{longtable}

%------------------------------------------------------------------------------
\section{Design patterns retenus}
%------------------------------------------------------------------------------
Le tableau suivant recense les patrons appliques, leur categorie (creationnel,
structurel, comportemental) et la couche concernee, avec la justification
\emph{avant / apres}.

\begin{longtable}{>{\bfseries}p{2.7cm} p{5.9cm} p{6.2cm}}
\toprule
\textbf{Pattern} & \textbf{Avant} & \textbf{Apres} \\
\midrule
\endhead
Composite \newline {\small\itshape structurel -- metier} &
Le calcul du miel parcourt manuellement les collections de cadres avec du code
de sommation duplique. &
\texttt{ElementRuche.getPoidsMiel()} est resolu recursivement
(Ruche \textrightarrow{} Compartiment \textrightarrow{} Cadre) : traitement
uniforme de l'arborescence. \\
State \newline {\small\itshape comportemental -- metier} &
Le statut de la ruche serait teste par des \texttt{switch} disperses (creee,
active, en collecte, cloturee, \ldots). &
\texttt{EtatRuche} et ses etats concrets encapsulent les transitions autorisees ;
alignement direct sur le diagramme de machine a etats. \\
Strategy \newline {\small\itshape comportemental -- metier} &
Les regles d'alerte, la conversion d'unites et le ROI seraient codes en dur. &
Familles d'algorithmes interchangeables (\texttt{RegleAlerte},
\texttt{ConvertisseurUnite}), injectables et testables isolement. \\
Observer \newline {\small\itshape comportemental -- metier} &
Le code de mesure appellerait directement les tableaux de bord (couplage fort). &
\texttt{GestionnaireAlertes} notifie les \texttt{Observateur} abonnes (tableau
de bord d'alertes, notification du responsable) au franchissement d'un seuil. \\
Command \newline {\small\itshape comportemental -- controle} &
Les operations terrain (collecte, division, changement de reine) sont des appels
directs, non rejouables hors-ligne. &
Chaque operation devient une \texttt{CommandeVisite} ; la
\texttt{FileCommandesHorsLigne} rejoue les commandes a la reconnexion
(mode hors-ligne + synchronisation differee). \\
Factory Method \newline {\small\itshape creationnel -- metier} &
Les ruches sont creees par des \texttt{new} disperses avec initialisation
manuelle des compartiments. &
\texttt{FabriqueRuche.creer(modele)} produit une ruche preconfiguree selon le
modele (Dadant, Langstroth\ldots) avec corps et capacites coherents. \\
Builder \newline {\small\itshape creationnel -- presentation} &
Le rapport de visite, aux nombreux champs optionnels, imposerait un constructeur
telescopique ou des setters epars. &
\texttt{ConstructeurRapport} construit \texttt{RapportVisite} de facon fluide
(constatations, actions, photos, evaluations). \\
Adapter \newline {\small\itshape structurel -- donnees} &
Les capteurs heterogenes (formats et unites differents) seraient traites au cas
par cas. &
\texttt{AdaptateurCapteur} uniformise chaque source en \texttt{Mesure} avec
unites parametrables (volet automatise). \\
Facade \newline {\small\itshape structurel -- controle} &
Le Servlet orchestre directement tous les services metier. &
\texttt{FacadeApplication} offre un point d'entree simplifie aux Servlets
par-dessus les sous-systemes. \\
Singleton \newline {\small\itshape creationnel -- transverse} &
Le fichier \texttt{ConfigBKS.ini} et les libelles i18n seraient relus en
plusieurs instances. &
\texttt{Configuration} et \texttt{GestionnaireI18n} garantissent une instance
unique partagee. \\
Proxy \newline {\small\itshape structurel -- services web} &
L'acces a l'API exposerait directement le service metier. &
Un \texttt{ProxySecurite} controle l'authentification OAuth et le chiffrement
SSL / X.509 avant de deleguer au service. \\
\bottomrule
\end{longtable}

%------------------------------------------------------------------------------
\section{Vue de conception refactoree}
%------------------------------------------------------------------------------
La figure~\ref{fig:solid} synthetise les principaux patrons dans un diagramme
de classes de conception : le Composite (composition de la ruche et calcul du
miel), le State (cycle de vie), les services segreges (ISP), la persistance par
abstraction (DIP / Repository), le couple Strategy + Observer (alertes) et le
trio Factory / Builder / Command.

\begin{figure}[htbp]
\centering
\includegraphics[width=\linewidth]{diagramme_classes_solid}
\caption{Diagramme de classes refactore : application de SOLID et des design
patterns au systeme Z\"umm.}
\label{fig:solid}
\end{figure}

\begin{encadre}
\textbf{Correspondance avec les couches :} les \texttt{Repository} appartiennent
a la couche d'acces aux donnees (JDBC), les \texttt{Service} et le domaine
(Composite, State, Strategy, Observer, Command) a la couche metier (EJB), la
\texttt{Facade} et le \texttt{Proxy} a la couche de controle (Servlets), le
\texttt{Builder} de rapport a la presentation (JSP). L'architecture en couches
faiblement couplees du chapitre~\ref{chap:conception} est ainsi renforcee, sans
etre modifiee.
\end{encadre}
 dans cahier_des_charges.tex
%  (utilise la charte "Miel & noir" : miel, mielfonce, encadre, colonne C,
%   longtable, booktabs -- deja definis dans le preambule principal)
%==============================================================================

%==============================================================================
\chapter{Annexe - Principes SOLID et design patterns}
\label{chap:solid}
%==============================================================================
Cette annexe repond a la question du \emph{comment} de la qualite logicielle.
Elle documente l'application des cinq principes \textbf{SOLID} et d'un ensemble
de \textbf{design patterns} (patrons de conception) au systeme \textbf{Z\"umm},
en mettant chaque decision en regard de la conception initiale (\emph{avant})
et de la conception refactoree (\emph{apres}). L'objectif est une architecture
\textbf{multi-niveaux, evolutive et faiblement couplee}, conforme aux exigences
du chapitre~\ref{chap:conception}.

\begin{encadre}
\textbf{Sources methodologiques :} les definitions retenues suivent la
presentation des principes SOLID de \emph{alexsoyes.com} et le catalogue de
patrons de \emph{refactoring.guru}, adaptes au domaine apicole du projet.
\end{encadre}

%------------------------------------------------------------------------------
\section{Les cinq principes SOLID}
%------------------------------------------------------------------------------
\begin{longtable}{>{\bfseries}p{2.4cm} p{6cm} p{6.2cm}}
\toprule
\textbf{Principe} & \textbf{Avant (conception initiale)} & \textbf{Apres (refactoree)} \\
\midrule
\endhead
S -- Responsabilite unique &
L'entite \texttt{Ruche} porte la methode de calcul
\texttt{getBeeHouseHoneyActualQuantity()} : elle melange l'etat (donnees) et la
logique metier ; \texttt{Visite} regroupe donnees et champs de rapport. &
Le calcul est extrait dans \texttt{ServiceMielImpl} (via le Composite) et le
rapport devient \texttt{RapportVisite}. Chaque classe n'a plus qu'une seule
raison de changer. \\
O -- Ouvert / ferme &
La detection d'alertes et la conversion d'unites reposeraient sur des
\texttt{if / switch} sur \texttt{typeSeuil} et \texttt{unite} : ajouter une
regle ou une unite impose de modifier le code existant. &
\texttt{RegleAlerte} (Strategy) et \texttt{ConvertisseurUnite} : une nouvelle
regle ou unite = une nouvelle classe, sans toucher au \texttt{GestionnaireAlertes}.
Ouvert a l'extension, ferme a la modification. \\
L -- Substitution de Liskov &
Le diagramme de cas d'utilisation generalise l'agent superviseur a partir de
l'agent apiculteur ; transpose tel quel en heritage de classes, un sous-type
restrictif casserait la substituabilite. &
\texttt{Corps} et \texttt{Hausse} heritent de \texttt{Compartiment} en
respectant integralement son contrat (capacite 1..10, \texttt{getPoidsMiel}).
Le role de l'agent passe par l'enumeration \texttt{Role} + permissions
(composition) plutot qu'un heritage restrictif. \\
I -- Segregation des interfaces &
Un service unique \og fourre-tout \fg{} (CRUD + visites + tableaux de bord +
API miel) obligerait l'application tierce a dependre de methodes inutilisees. &
Interfaces ciblees : \texttt{ServiceRequeteMiel} (l'API tierce ne voit que
\texttt{getBeeHouseHoneyActualQuantity}), \texttt{ServiceVisite},
\texttt{ServiceTableauDeBord}. Chaque client ne depend que de son contrat. \\
D -- Inversion de dependance &
La couche metier (EJB) appelle directement JDBC (implementation concrete) :
couplage fort a la base et aux \texttt{ResultSet}. &
La couche metier depend d'abstractions \texttt{RucheRepository} /
\texttt{VisiteRepository} (pattern DAO), implementees par
\texttt{RucheRepositoryJdbc}. Tests par mock et migration JDBC \textrightarrow{}
Spring Data JPA facilites (cf. annexe~\ref{chap:migration}). \\
\bottomrule
\end{longtable}

%------------------------------------------------------------------------------
\section{Design patterns retenus}
%------------------------------------------------------------------------------
Le tableau suivant recense les patrons appliques, leur categorie (creationnel,
structurel, comportemental) et la couche concernee, avec la justification
\emph{avant / apres}.

\begin{longtable}{>{\bfseries}p{2.7cm} p{5.9cm} p{6.2cm}}
\toprule
\textbf{Pattern} & \textbf{Avant} & \textbf{Apres} \\
\midrule
\endhead
Composite \newline {\small\itshape structurel -- metier} &
Le calcul du miel parcourt manuellement les collections de cadres avec du code
de sommation duplique. &
\texttt{ElementRuche.getPoidsMiel()} est resolu recursivement
(Ruche \textrightarrow{} Compartiment \textrightarrow{} Cadre) : traitement
uniforme de l'arborescence. \\
State \newline {\small\itshape comportemental -- metier} &
Le statut de la ruche serait teste par des \texttt{switch} disperses (creee,
active, en collecte, cloturee, \ldots). &
\texttt{EtatRuche} et ses etats concrets encapsulent les transitions autorisees ;
alignement direct sur le diagramme de machine a etats. \\
Strategy \newline {\small\itshape comportemental -- metier} &
Les regles d'alerte, la conversion d'unites et le ROI seraient codes en dur. &
Familles d'algorithmes interchangeables (\texttt{RegleAlerte},
\texttt{ConvertisseurUnite}), injectables et testables isolement. \\
Observer \newline {\small\itshape comportemental -- metier} &
Le code de mesure appellerait directement les tableaux de bord (couplage fort). &
\texttt{GestionnaireAlertes} notifie les \texttt{Observateur} abonnes (tableau
de bord d'alertes, notification du responsable) au franchissement d'un seuil. \\
Command \newline {\small\itshape comportemental -- controle} &
Les operations terrain (collecte, division, changement de reine) sont des appels
directs, non rejouables hors-ligne. &
Chaque operation devient une \texttt{CommandeVisite} ; la
\texttt{FileCommandesHorsLigne} rejoue les commandes a la reconnexion
(mode hors-ligne + synchronisation differee). \\
Factory Method \newline {\small\itshape creationnel -- metier} &
Les ruches sont creees par des \texttt{new} disperses avec initialisation
manuelle des compartiments. &
\texttt{FabriqueRuche.creer(modele)} produit une ruche preconfiguree selon le
modele (Dadant, Langstroth\ldots) avec corps et capacites coherents. \\
Builder \newline {\small\itshape creationnel -- presentation} &
Le rapport de visite, aux nombreux champs optionnels, imposerait un constructeur
telescopique ou des setters epars. &
\texttt{ConstructeurRapport} construit \texttt{RapportVisite} de facon fluide
(constatations, actions, photos, evaluations). \\
Adapter \newline {\small\itshape structurel -- donnees} &
Les capteurs heterogenes (formats et unites differents) seraient traites au cas
par cas. &
\texttt{AdaptateurCapteur} uniformise chaque source en \texttt{Mesure} avec
unites parametrables (volet automatise). \\
Facade \newline {\small\itshape structurel -- controle} &
Le Servlet orchestre directement tous les services metier. &
\texttt{FacadeApplication} offre un point d'entree simplifie aux Servlets
par-dessus les sous-systemes. \\
Singleton \newline {\small\itshape creationnel -- transverse} &
Le fichier \texttt{ConfigBKS.ini} et les libelles i18n seraient relus en
plusieurs instances. &
\texttt{Configuration} et \texttt{GestionnaireI18n} garantissent une instance
unique partagee. \\
Proxy \newline {\small\itshape structurel -- services web} &
L'acces a l'API exposerait directement le service metier. &
Un \texttt{ProxySecurite} controle l'authentification OAuth et le chiffrement
SSL / X.509 avant de deleguer au service. \\
\bottomrule
\end{longtable}

%------------------------------------------------------------------------------
\section{Vue de conception refactoree}
%------------------------------------------------------------------------------
La figure~\ref{fig:solid} synthetise les principaux patrons dans un diagramme
de classes de conception : le Composite (composition de la ruche et calcul du
miel), le State (cycle de vie), les services segreges (ISP), la persistance par
abstraction (DIP / Repository), le couple Strategy + Observer (alertes) et le
trio Factory / Builder / Command.

\begin{figure}[htbp]
\centering
\includegraphics[width=\linewidth]{diagramme_classes_solid}
\caption{Diagramme de classes refactore : application de SOLID et des design
patterns au systeme Z\"umm.}
\label{fig:solid}
\end{figure}

\begin{encadre}
\textbf{Correspondance avec les couches :} les \texttt{Repository} appartiennent
a la couche d'acces aux donnees (JDBC), les \texttt{Service} et le domaine
(Composite, State, Strategy, Observer, Command) a la couche metier (EJB), la
\texttt{Facade} et le \texttt{Proxy} a la couche de controle (Servlets), le
\texttt{Builder} de rapport a la presentation (JSP). L'architecture en couches
faiblement couplees du chapitre~\ref{chap:conception} est ainsi renforcee, sans
etre modifiee.
\end{encadre}
 dans cahier_des_charges.tex
%  (utilise la charte "Miel & noir" : miel, mielfonce, encadre, colonne C,
%   longtable, booktabs -- deja definis dans le preambule principal)
%==============================================================================

%==============================================================================
\chapter{Annexe - Principes SOLID et design patterns}
\label{chap:solid}
%==============================================================================
Cette annexe repond a la question du \emph{comment} de la qualite logicielle.
Elle documente l'application des cinq principes \textbf{SOLID} et d'un ensemble
de \textbf{design patterns} (patrons de conception) au systeme \textbf{Z\"umm},
en mettant chaque decision en regard de la conception initiale (\emph{avant})
et de la conception refactoree (\emph{apres}). L'objectif est une architecture
\textbf{multi-niveaux, evolutive et faiblement couplee}, conforme aux exigences
du chapitre~\ref{chap:conception}.

\begin{encadre}
\textbf{Sources methodologiques :} les definitions retenues suivent la
presentation des principes SOLID de \emph{alexsoyes.com} et le catalogue de
patrons de \emph{refactoring.guru}, adaptes au domaine apicole du projet.
\end{encadre}

%------------------------------------------------------------------------------
\section{Les cinq principes SOLID}
%------------------------------------------------------------------------------
\begin{longtable}{>{\bfseries}p{2.4cm} p{6cm} p{6.2cm}}
\toprule
\textbf{Principe} & \textbf{Avant (conception initiale)} & \textbf{Apres (refactoree)} \\
\midrule
\endhead
S -- Responsabilite unique &
L'entite \texttt{Ruche} porte la methode de calcul
\texttt{getBeeHouseHoneyActualQuantity()} : elle melange l'etat (donnees) et la
logique metier ; \texttt{Visite} regroupe donnees et champs de rapport. &
Le calcul est extrait dans \texttt{ServiceMielImpl} (via le Composite) et le
rapport devient \texttt{RapportVisite}. Chaque classe n'a plus qu'une seule
raison de changer. \\
O -- Ouvert / ferme &
La detection d'alertes et la conversion d'unites reposeraient sur des
\texttt{if / switch} sur \texttt{typeSeuil} et \texttt{unite} : ajouter une
regle ou une unite impose de modifier le code existant. &
\texttt{RegleAlerte} (Strategy) et \texttt{ConvertisseurUnite} : une nouvelle
regle ou unite = une nouvelle classe, sans toucher au \texttt{GestionnaireAlertes}.
Ouvert a l'extension, ferme a la modification. \\
L -- Substitution de Liskov &
Le diagramme de cas d'utilisation generalise l'agent superviseur a partir de
l'agent apiculteur ; transpose tel quel en heritage de classes, un sous-type
restrictif casserait la substituabilite. &
\texttt{Corps} et \texttt{Hausse} heritent de \texttt{Compartiment} en
respectant integralement son contrat (capacite 1..10, \texttt{getPoidsMiel}).
Le role de l'agent passe par l'enumeration \texttt{Role} + permissions
(composition) plutot qu'un heritage restrictif. \\
I -- Segregation des interfaces &
Un service unique \og fourre-tout \fg{} (CRUD + visites + tableaux de bord +
API miel) obligerait l'application tierce a dependre de methodes inutilisees. &
Interfaces ciblees : \texttt{ServiceRequeteMiel} (l'API tierce ne voit que
\texttt{getBeeHouseHoneyActualQuantity}), \texttt{ServiceVisite},
\texttt{ServiceTableauDeBord}. Chaque client ne depend que de son contrat. \\
D -- Inversion de dependance &
La couche metier (EJB) appelle directement JDBC (implementation concrete) :
couplage fort a la base et aux \texttt{ResultSet}. &
La couche metier depend d'abstractions \texttt{RucheRepository} /
\texttt{VisiteRepository} (pattern DAO), implementees par
\texttt{RucheRepositoryJdbc}. Tests par mock et migration JDBC \textrightarrow{}
Spring Data JPA facilites (cf. annexe~\ref{chap:migration}). \\
\bottomrule
\end{longtable}

%------------------------------------------------------------------------------
\section{Design patterns retenus}
%------------------------------------------------------------------------------
Le tableau suivant recense les patrons appliques, leur categorie (creationnel,
structurel, comportemental) et la couche concernee, avec la justification
\emph{avant / apres}.

\begin{longtable}{>{\bfseries}p{2.7cm} p{5.9cm} p{6.2cm}}
\toprule
\textbf{Pattern} & \textbf{Avant} & \textbf{Apres} \\
\midrule
\endhead
Composite \newline {\small\itshape structurel -- metier} &
Le calcul du miel parcourt manuellement les collections de cadres avec du code
de sommation duplique. &
\texttt{ElementRuche.getPoidsMiel()} est resolu recursivement
(Ruche \textrightarrow{} Compartiment \textrightarrow{} Cadre) : traitement
uniforme de l'arborescence. \\
State \newline {\small\itshape comportemental -- metier} &
Le statut de la ruche serait teste par des \texttt{switch} disperses (creee,
active, en collecte, cloturee, \ldots). &
\texttt{EtatRuche} et ses etats concrets encapsulent les transitions autorisees ;
alignement direct sur le diagramme de machine a etats. \\
Strategy \newline {\small\itshape comportemental -- metier} &
Les regles d'alerte, la conversion d'unites et le ROI seraient codes en dur. &
Familles d'algorithmes interchangeables (\texttt{RegleAlerte},
\texttt{ConvertisseurUnite}), injectables et testables isolement. \\
Observer \newline {\small\itshape comportemental -- metier} &
Le code de mesure appellerait directement les tableaux de bord (couplage fort). &
\texttt{GestionnaireAlertes} notifie les \texttt{Observateur} abonnes (tableau
de bord d'alertes, notification du responsable) au franchissement d'un seuil. \\
Command \newline {\small\itshape comportemental -- controle} &
Les operations terrain (collecte, division, changement de reine) sont des appels
directs, non rejouables hors-ligne. &
Chaque operation devient une \texttt{CommandeVisite} ; la
\texttt{FileCommandesHorsLigne} rejoue les commandes a la reconnexion
(mode hors-ligne + synchronisation differee). \\
Factory Method \newline {\small\itshape creationnel -- metier} &
Les ruches sont creees par des \texttt{new} disperses avec initialisation
manuelle des compartiments. &
\texttt{FabriqueRuche.creer(modele)} produit une ruche preconfiguree selon le
modele (Dadant, Langstroth\ldots) avec corps et capacites coherents. \\
Builder \newline {\small\itshape creationnel -- presentation} &
Le rapport de visite, aux nombreux champs optionnels, imposerait un constructeur
telescopique ou des setters epars. &
\texttt{ConstructeurRapport} construit \texttt{RapportVisite} de facon fluide
(constatations, actions, photos, evaluations). \\
Adapter \newline {\small\itshape structurel -- donnees} &
Les capteurs heterogenes (formats et unites differents) seraient traites au cas
par cas. &
\texttt{AdaptateurCapteur} uniformise chaque source en \texttt{Mesure} avec
unites parametrables (volet automatise). \\
Facade \newline {\small\itshape structurel -- controle} &
Le Servlet orchestre directement tous les services metier. &
\texttt{FacadeApplication} offre un point d'entree simplifie aux Servlets
par-dessus les sous-systemes. \\
Singleton \newline {\small\itshape creationnel -- transverse} &
Le fichier \texttt{ConfigBKS.ini} et les libelles i18n seraient relus en
plusieurs instances. &
\texttt{Configuration} et \texttt{GestionnaireI18n} garantissent une instance
unique partagee. \\
Proxy \newline {\small\itshape structurel -- services web} &
L'acces a l'API exposerait directement le service metier. &
Un \texttt{ProxySecurite} controle l'authentification OAuth et le chiffrement
SSL / X.509 avant de deleguer au service. \\
\bottomrule
\end{longtable}

%------------------------------------------------------------------------------
\section{Vue de conception refactoree}
%------------------------------------------------------------------------------
La figure~\ref{fig:solid} synthetise les principaux patrons dans un diagramme
de classes de conception : le Composite (composition de la ruche et calcul du
miel), le State (cycle de vie), les services segreges (ISP), la persistance par
abstraction (DIP / Repository), le couple Strategy + Observer (alertes) et le
trio Factory / Builder / Command.

\begin{figure}[htbp]
\centering
\includegraphics[width=\linewidth]{diagramme_classes_solid}
\caption{Diagramme de classes refactore : application de SOLID et des design
patterns au systeme Z\"umm.}
\label{fig:solid}
\end{figure}

\begin{encadre}
\textbf{Correspondance avec les couches :} les \texttt{Repository} appartiennent
a la couche d'acces aux donnees (JDBC), les \texttt{Service} et le domaine
(Composite, State, Strategy, Observer, Command) a la couche metier (EJB), la
\texttt{Facade} et le \texttt{Proxy} a la couche de controle (Servlets), le
\texttt{Builder} de rapport a la presentation (JSP). L'architecture en couches
faiblement couplees du chapitre~\ref{chap:conception} est ainsi renforcee, sans
etre modifiee.
\end{encadre}
 dans cahier_des_charges.tex
%  (utilise la charte "Miel & noir" : miel, mielfonce, encadre, colonne C,
%   longtable, booktabs -- deja definis dans le preambule principal)
%==============================================================================

%==============================================================================
\chapter{Annexe - Principes SOLID et design patterns}
\label{chap:solid}
%==============================================================================
Cette annexe repond a la question du \emph{comment} de la qualite logicielle.
Elle documente l'application des cinq principes \textbf{SOLID} et d'un ensemble
de \textbf{design patterns} (patrons de conception) au systeme \textbf{Z\"umm},
en mettant chaque decision en regard de la conception initiale (\emph{avant})
et de la conception refactoree (\emph{apres}). L'objectif est une architecture
\textbf{multi-niveaux, evolutive et faiblement couplee}, conforme aux exigences
du chapitre~\ref{chap:conception}.

\begin{encadre}
\textbf{Sources methodologiques :} les definitions retenues suivent la
presentation des principes SOLID de \emph{alexsoyes.com} et le catalogue de
patrons de \emph{refactoring.guru}, adaptes au domaine apicole du projet.
\end{encadre}

%------------------------------------------------------------------------------
\section{Les cinq principes SOLID}
%------------------------------------------------------------------------------
\begin{longtable}{>{\bfseries}p{2.4cm} p{6cm} p{6.2cm}}
\toprule
\textbf{Principe} & \textbf{Avant (conception initiale)} & \textbf{Apres (refactoree)} \\
\midrule
\endhead
S -- Responsabilite unique &
L'entite \texttt{Ruche} porte la methode de calcul
\texttt{getBeeHouseHoneyActualQuantity()} : elle melange l'etat (donnees) et la
logique metier ; \texttt{Visite} regroupe donnees et champs de rapport. &
Le calcul est extrait dans \texttt{ServiceMielImpl} (via le Composite) et le
rapport devient \texttt{RapportVisite}. Chaque classe n'a plus qu'une seule
raison de changer. \\
O -- Ouvert / ferme &
La detection d'alertes et la conversion d'unites reposeraient sur des
\texttt{if / switch} sur \texttt{typeSeuil} et \texttt{unite} : ajouter une
regle ou une unite impose de modifier le code existant. &
\texttt{RegleAlerte} (Strategy) et \texttt{ConvertisseurUnite} : une nouvelle
regle ou unite = une nouvelle classe, sans toucher au \texttt{GestionnaireAlertes}.
Ouvert a l'extension, ferme a la modification. \\
L -- Substitution de Liskov &
Le diagramme de cas d'utilisation generalise l'agent superviseur a partir de
l'agent apiculteur ; transpose tel quel en heritage de classes, un sous-type
restrictif casserait la substituabilite. &
\texttt{Corps} et \texttt{Hausse} heritent de \texttt{Compartiment} en
respectant integralement son contrat (capacite 1..10, \texttt{getPoidsMiel}).
Le role de l'agent passe par l'enumeration \texttt{Role} + permissions
(composition) plutot qu'un heritage restrictif. \\
I -- Segregation des interfaces &
Un service unique \og fourre-tout \fg{} (CRUD + visites + tableaux de bord +
API miel) obligerait l'application tierce a dependre de methodes inutilisees. &
Interfaces ciblees : \texttt{ServiceRequeteMiel} (l'API tierce ne voit que
\texttt{getBeeHouseHoneyActualQuantity}), \texttt{ServiceVisite},
\texttt{ServiceTableauDeBord}. Chaque client ne depend que de son contrat. \\
D -- Inversion de dependance &
La couche metier (EJB) appelle directement JDBC (implementation concrete) :
couplage fort a la base et aux \texttt{ResultSet}. &
La couche metier depend d'abstractions \texttt{RucheRepository} /
\texttt{VisiteRepository} (pattern DAO), implementees par
\texttt{RucheRepositoryJdbc}. Tests par mock et migration JDBC \textrightarrow{}
Spring Data JPA facilites (cf. annexe~\ref{chap:migration}). \\
\bottomrule
\end{longtable}

%------------------------------------------------------------------------------
\section{Design patterns retenus}
%------------------------------------------------------------------------------
Le tableau suivant recense les patrons appliques, leur categorie (creationnel,
structurel, comportemental) et la couche concernee, avec la justification
\emph{avant / apres}.

\begin{longtable}{>{\bfseries}p{2.7cm} p{5.9cm} p{6.2cm}}
\toprule
\textbf{Pattern} & \textbf{Avant} & \textbf{Apres} \\
\midrule
\endhead
Composite \newline {\small\itshape structurel -- metier} &
Le calcul du miel parcourt manuellement les collections de cadres avec du code
de sommation duplique. &
\texttt{ElementRuche.getPoidsMiel()} est resolu recursivement
(Ruche \textrightarrow{} Compartiment \textrightarrow{} Cadre) : traitement
uniforme de l'arborescence. \\
State \newline {\small\itshape comportemental -- metier} &
Le statut de la ruche serait teste par des \texttt{switch} disperses (creee,
active, en collecte, cloturee, \ldots). &
\texttt{EtatRuche} et ses etats concrets encapsulent les transitions autorisees ;
alignement direct sur le diagramme de machine a etats. \\
Strategy \newline {\small\itshape comportemental -- metier} &
Les regles d'alerte, la conversion d'unites et le ROI seraient codes en dur. &
Familles d'algorithmes interchangeables (\texttt{RegleAlerte},
\texttt{ConvertisseurUnite}), injectables et testables isolement. \\
Observer \newline {\small\itshape comportemental -- metier} &
Le code de mesure appellerait directement les tableaux de bord (couplage fort). &
\texttt{GestionnaireAlertes} notifie les \texttt{Observateur} abonnes (tableau
de bord d'alertes, notification du responsable) au franchissement d'un seuil. \\
Command \newline {\small\itshape comportemental -- controle} &
Les operations terrain (collecte, division, changement de reine) sont des appels
directs, non rejouables hors-ligne. &
Chaque operation devient une \texttt{CommandeVisite} ; la
\texttt{FileCommandesHorsLigne} rejoue les commandes a la reconnexion
(mode hors-ligne + synchronisation differee). \\
Factory Method \newline {\small\itshape creationnel -- metier} &
Les ruches sont creees par des \texttt{new} disperses avec initialisation
manuelle des compartiments. &
\texttt{FabriqueRuche.creer(modele)} produit une ruche preconfiguree selon le
modele (Dadant, Langstroth\ldots) avec corps et capacites coherents. \\
Builder \newline {\small\itshape creationnel -- presentation} &
Le rapport de visite, aux nombreux champs optionnels, imposerait un constructeur
telescopique ou des setters epars. &
\texttt{ConstructeurRapport} construit \texttt{RapportVisite} de facon fluide
(constatations, actions, photos, evaluations). \\
Adapter \newline {\small\itshape structurel -- donnees} &
Les capteurs heterogenes (formats et unites differents) seraient traites au cas
par cas. &
\texttt{AdaptateurCapteur} uniformise chaque source en \texttt{Mesure} avec
unites parametrables (volet automatise). \\
Facade \newline {\small\itshape structurel -- controle} &
Le Servlet orchestre directement tous les services metier. &
\texttt{FacadeApplication} offre un point d'entree simplifie aux Servlets
par-dessus les sous-systemes. \\
Singleton \newline {\small\itshape creationnel -- transverse} &
Le fichier \texttt{ConfigBKS.ini} et les libelles i18n seraient relus en
plusieurs instances. &
\texttt{Configuration} et \texttt{GestionnaireI18n} garantissent une instance
unique partagee. \\
Proxy \newline {\small\itshape structurel -- services web} &
L'acces a l'API exposerait directement le service metier. &
Un \texttt{ProxySecurite} controle l'authentification OAuth et le chiffrement
SSL / X.509 avant de deleguer au service. \\
\bottomrule
\end{longtable}

%------------------------------------------------------------------------------
\section{Vue de conception refactoree}
%------------------------------------------------------------------------------
La figure~\ref{fig:solid} synthetise les principaux patrons dans un diagramme
de classes de conception : le Composite (composition de la ruche et calcul du
miel), le State (cycle de vie), les services segreges (ISP), la persistance par
abstraction (DIP / Repository), le couple Strategy + Observer (alertes) et le
trio Factory / Builder / Command.

\begin{figure}[htbp]
\centering
\includegraphics[width=\linewidth]{diagramme_classes_solid}
\caption{Diagramme de classes refactore : application de SOLID et des design
patterns au systeme Z\"umm.}
\label{fig:solid}
\end{figure}

\begin{encadre}
\textbf{Correspondance avec les couches :} les \texttt{Repository} appartiennent
a la couche d'acces aux donnees (JDBC), les \texttt{Service} et le domaine
(Composite, State, Strategy, Observer, Command) a la couche metier (EJB), la
\texttt{Facade} et le \texttt{Proxy} a la couche de controle (Servlets), le
\texttt{Builder} de rapport a la presentation (JSP). L'architecture en couches
faiblement couplees du chapitre~\ref{chap:conception} est ainsi renforcee, sans
etre modifiee.
\end{encadre}
